
\svnid{$Id: $}
\svnid{$URL: $}

\section{FOR REVIEW: Junction advection for steady state conditions (flow diversion and deceleration at junction; non equidistant grid) applying conservation of momentum (iadvec1d=1)}
\newrefsegment

%---------------------------------------------------------------------------------
\section*{Quality Assurance}
\begin{tabular}{@{}p{27mm}@{}p{0mm}p{\textwidth-52mm-48pt}p{15mm}p{10mm}}
\textbf{Date} & \textbf{:} &   May 13 2019\\ 
\textbf{Author} & \textbf{:} & Christianne Luger  & \textbf{Initials:} & CL \\
\textbf{Review} & \textbf{:} &   & \textbf{Initials:} & \ldots 
\end{tabular}

%---------------------------------------------------------------------------------
\section*{Version information}
{{
\begin{tabular}{@{}p{27mm}@{}p{0mm}p{\textwidth-27mm-24pt}}
\textbf{Executable}   & \textbf{:} & Deltares, D-Flow FM Version 1.2.54.64224M, Jun 28 2019, 15:14:31  \\
\textbf{Location input}     & \textbf{:} & \url{https://repos.deltares.nl/repos/DSCTestbench/trunk/cases/e106_dflow1d/f04_numerical-aspects/c10_junction-advection04_iadvec1d_1} \\
\textbf{Revision input} & \textbf{:} & - \\
\textbf{Location output}     & \textbf{:} & \url{https://repos.deltares.nl/repos/DSCTestbench/trunk/cases/e02_dflowfm/f104_1D_numerical-aspects/c04_junction-advection-deceleration-non-equidistant} \\
\textbf{Revision output} & \textbf{:} & -
\end{tabular}
}}

%---------------------------------------------------------------------------------
\section*{Purpose}
\textbf{FOR REVIEW}
The purpose of test \textbf{e02}(dflowfm)\_\textbf{f104}(numerical-aspects)\_\textbf{c04}(junction-advection-deceleration-non-equidistant) is to verify the implementation of junction advection (see \autoref{10404:description of junction advection}) for steady state conditions (flow diversion and deceleration at junction; non equidistant grid) while applying conservation of momentum (\textbf{\textcolor[rgb]{1,0,0}{iadvec1d=1}}).

Another objective of the test is to verify that the constant discharge and constant water level boundaries work as expected.

%---------------------------------------------------------------------------------
\section*{Linked claims}
\begin{enumerate}
\item Junction advection is correctly implemented for steady state conditions (flow diversion and deceleration at junction; non equidistant grid) while applying conservation of momentum (iadvec1d=1). 
\item Constant (upstream) discharge boundaries are correctly applied on the systems.
\item Constant (downstream) water level boundaries are correctly applied on the systems.
\end{enumerate}

%---------------------------------------------------------------------------------
\section*{Approach}

Water levels and velocities near junctions are compared with the results of  models having no junctions (or connection nodes). More precisely, results of sub-models \textbf{\textcolor[rgb]{1,0,0}{T2}} to \textbf{\textcolor[rgb]{1,0,0}{T4}} are validated against the results of sub-model \textbf{\textcolor[rgb]{1,0,0}{T1}} (see \autoref{fig:10404:Test model lay-out}). The results are further compared to a SOBEK3 identical test case  \textbf{e106}(dflow1d)\_\textbf{f04}(numerical-aspects)\_\textbf{c10}(junction-advection04-iadvec1d=1).

\begin{figure}[H]
    \centering
    \includegraphics*[width=0.95\textwidth]{figures/e02_f104_c04_lay-out.jpg}
    \caption{Lay-out of the test model}
    \label{fig:10404:Test model lay-out}
\end{figure} 

%---------------------------------------------------------------------------------
\section*{Conclusion} 

A comparison of the results of this case ran on DFlowFM to an identical test case ran on SOBEK3 show slight differences in water levels, water depths, discharges and velocities at and around the connection nodes. However, these differences are not due to the incorrect application of junction advection and do not cause any additional numerical disturbances. Therefore, the implementation of junction advection for steady state conditions  (flow diversion and deceleration at junction; non-equidistant grid) while applying conservation of momentum (iadvec1d=1) is accepted. The implementation of a constant upstream discharge boundary and constant downstream water level boundary is also accepted.

%---------------------------------------------------------------------------------
\section*{Model description}

The test comprises of four separate sub-models (see \autoref{fig:10404:Test model lay-out}), having a \underline{non-equidistant} grid-spacing.
\begin{enumerate}
	\item Sub-model \textbf{\textcolor[rgb]{1,0,0}{T1}} comprises of one branch \textbf{B1} (length 25000 m) only (see \autoref{fig:10404:Test model lay-out}).
		\begin{itemize}
			\item At T1\_Bnd\_B1\_x0m a \textbf{constant} discharge of 790.57 $m^3$/s is imposed. 
			\item At T1\_Bnd\_B1\_x25000m a \textbf{constant} water level of 2.5 m is imposed.
			\item Branch B1 has a non-equidistant grid. Water-level-points are located at x=0m, x=2500m, x=5000m, x=7500m, x=10000m, x=11000m, x=12500m, x=14500m, x=15000m, x=16500m, x=17500m, x=19000m, x=20000m,x=22500m and x=25000m.
			\item Branch B1 contains three "open Y-Z rectangular cross-sections" (width=\textcolor[rgb]{0.071,0,1}{\textcolor[rgb]{0.5,0,0.15}{\textbf{500}}} m, Chezy=50 $m^{0.5}$/s), respectively having an elevation of \textbf{15.10} m AD at T1\_Bnd\_B1\_x0m, \textbf{0.10} m AD at fixed water-level-point T1\_h\_pnt\_x15000m, and \textbf{0.00} m AD at \newline T1\_Bnd\_B1\_x25000m. 
			\item The flow decelerates locally at the fixed water-level-point T1\_h\_pnt\_x15000m from 0.80 to 0.55 m/s.
		\end{itemize} \


	\item Sub-model \textbf{\textcolor[rgb]{1,0,0}{T2}} comprises of two branches \textbf{B1} (length 15000 m) and \textbf{B2} (length 10000 m), that are mutually connected at connection node T2\_ConNode (see \autoref{fig:10404:Test model lay-out}). 
		\begin{itemize}
			\item At T2\_Bnd\_B1\_x0m a \textbf{constant} discharge of 790.57 $m^3$/s is imposed.
			\item At T2\_Bnd\_B2\_x10000m a \textbf{constant} water level of 2.5 m is imposed.
			\item Branch B1 has a non-equidistant grid. Water-level-points are located at x=0m, x=2500m, x=5000m, x=7500m, x=10000m, x=11000m, x=12500m, x=14500m and x=15000m.
			\item Branch B2 has a non-equidistant grid. Water-level-points are located at x=0m, x=1500m, x=2500m, x=4000m, x=5000m, x=7500m and x=10000m.
			\item Branch B1 contains two "open Y-Z rectangular cross-sections" (width=\textbf{\textcolor[rgb]{0.5,0,0.15}{500}} m, Chezy=50 $m^{0.5}$/s), respectively having an elevation of \textbf{15.10} m AD at T2\_Bnd\_B1\_x0m and \textbf{0.10} m AD at T2\_ConNode.
			\item Branch B2 contains two "open Y-Z rectangular cross-sections" (width=\textbf{\textcolor[rgb]{0.5,0,0.15}{500}} m, Chezy=50 $m^{0.5}$/s), respectively having an elevation \textbf{0.10} m at T2\_ConNode and \textbf{0.00} m AD at \newline T2\_Bnd\_B2\_x10000m.
		\end{itemize} 
		
		\Note The bathymetry and hydraulic properties of branch B1 and B2 of sub-model \textbf{\textcolor[rgb]{1,0,0}{T2}} are indentical to those branch B1 of sub-model \textbf{\textcolor[rgb]{1,0,0}{T1}}. \\
	
	\item Sub-model \textbf{\textcolor[rgb]{1,0,0}{T3}} comprises of three branches \textbf{B1} (length 15000 m), \textbf{B2} (length 10000 m) and \textbf{B3} (length 10000 m), that are mutually connected at connection node T3\_ConNode (see \autoref{fig:10404:Test model lay-out}).
		\begin{itemize}
			\item At T3\_Bnd\_B1\_x0m a \textbf{constant} discharge of 790.57 $m^3$/s is imposed.
			\item At T3\_Bnd\_B2\_x10000m a \textbf{constant} water level of 2.5 m is imposed.
			\item At T3\_Bnd\_B3\_x10000m a \textbf{constant} water level of 2.5 m is imposed.		
			\item Branch B1 has a non-equidistant grid. Water-level-points are located at x=0m, x=2500m, x=5000m, x=7500m, x=10000m, x=11000m, x=12500m, x=14500m and x=15000m.
			\item Branch B2 has a non-equidistant grid. Water-level-points are located at x=0m, x=1500m, x=2500m, x=4000m, x=5000m, x=7500m and x=10000m.
			\item Branch B3 has a non-equidistant grid. Water-level-points are located at x=0m, x=1500m, x=2500m, x=4000m, x=5000m, x=7500m and x=10000m.
			\item Branch B1 contains two "open Y-Z rectangular cross-sections" (width=\textbf{\textcolor[rgb]{0.5,0,0.15}{500}} m, Chezy=50 $m^{0.5}$/s), respectively having an elevation of \textbf{15.1} m AD at T3\_Bnd\_B1\_x0m and \textbf{0.10} m AD at T3\_ConNode. 
			\item Branch B2 contains two "open Y-Z rectangular cross-sections" (width=\textbf{\textcolor[rgb]{0.5,0,0.15}{250}} m, Chezy=50 $m^{0.5}$/s), respectively having an elevation of \textbf{0.10} m at T3\_ConNode and \textbf{0.00} m AD at \newline T3\_Bnd\_B2\_x10000m.
			\item Branch B3 contains two "open Y-Z rectangular cross-sections" (width=\textbf{\textcolor[rgb]{0.5,0,0.15}{250}} m, Chezy=50 $m^{0.5}$/s), respectively having an elevation of \textbf{0.10} m at T3\_ConNode and \textbf{0.00} m AD at \newline T3\_Bnd\_B3\_x10000m.
		\end{itemize} 
			
		\Note The bathymetry and hydraulic properties per meter of cross-sectional width of branch B1 and B2 of sub-model \textbf{\textcolor[rgb]{1,0,0}{T3}} are indentical to those of branch B1 of sub-model \textbf{\textcolor[rgb]{1,0,0}{T1}}. The same applies for branch B1 and B3 of sub-model \textbf{\textcolor[rgb]{1,0,0}{T3}}.\\
	
	\item Sub-model \textbf{\textcolor[rgb]{1,0,0}{T4}} comprises of four branches \textbf{B1} (length 15000 m), \textbf{B2} (length 10000 m), \textbf{B3} (length 10000 m), and \textbf{B4} (length 15000 m), that are mutually connected at connection node T4\_ConNode (see \autoref{fig:10404:Test model lay-out}).
		\begin{itemize}
			\item At T4\_Bnd\_B1\_x0m a \textbf{constant} discharge of 237.171 $m^3$/s is imposed.
			\item At T4\_Bnd\_B2\_x10000m a \textbf{constant} water level of 2.5 m is imposed.
			\item At T4\_Bnd\_B3\_x10000m a \textbf{constant} water level of 2.5 m is imposed.
			\item At T4\_Bnd\_B4\_x0m a \textbf{constant} discharge of 553.399 $m^3$/s is imposed.
			\item Branch B1 has a non-equidistant grid. Water-level-points are located at x=0m, x=2500m, x=5000m, x=7500m, x=10000m, x=11000m, x=12500m, x=14500m and x=15000m.
			\item Branch B2 has a non-equidistant grid. Water-level-points are located at x=0m, x=1500m, x=2500m, x=4000m, x=5000m, x=7500m and x=10000m.
			\item Branch B3 has a non-equidistant grid. Water-level-points are located at x=0m, x=1500m, x=2500m, x=4000m, x=5000m, x=7500m and x=10000m.
			\item Branch B4 has a non-equidistant grid. Water-level-points are located at x=0m, x=2500m, x=5000m, x=7500m, x=10000m, x=11000m, x=12500m, x=14500m and x=15000m.
			\item Branch B1 contains two "open Y-Z rectangular cross-sections" (width=\textbf{\textcolor[rgb]{0.5,0,0.15}{150}} m, Chezy=50 $m^{0.5}$/s), respectively having an elevation of \textbf{15.10} m AD at T4\_Bnd\_B1\_x0m and \textbf{0.10} m AD at T4\_ConNode. 
			\item Branch B2 contains two "open Y-Z rectangular cross-sections" (width=\textbf{\textcolor[rgb]{0.5,0,0.15}{50}} m, Chezy=50 $m^{0.5}$/s), respectively having an elevation \textbf{0.10} m AD at T4\_ConNode and \textbf{0.00} m AD at \newline T4\_Bnd\_B2\_x10000m.
			\item Branch B3 contains two "open Y-Z rectangular cross-sections" (width=\textbf{\textcolor[rgb]{0.5,0,0.15}{450}} m, Chezy=50 $m^{0.5}$/s), respectively having an elevation \textbf{0.10} m AD at T4\_ConNode and \textbf{0.00} m AD at T4\_Bnd\_B3\_x10000m.
			\item Branch B4 contains two "open Y-Z rectangular cross-sections" (width=\textbf{\textcolor[rgb]{0.5,0,0.15}{350}} m, Chezy=50 $m^{0.5}$/s), respectively having an elevation of \textbf{15.10} m AD at T4\_Bnd\_B4\_x0m and \textbf{10.00} m AD at T4\_ConNode.
		\end{itemize}
		
				\Note The bathymetry and hydraulic properties per meter of cross-sectional width of branch B1 and B2 of sub-model \textbf{\textcolor[rgb]{1,0,0}{T4}} are indentical to those of branch B1 of sub-model \textbf{\textcolor[rgb]{1,0,0}{T1}}. The same applies for branch B1 and B3, for branch B4 and B2, and for branch B4 and B3 of sub-model \textbf{\textcolor[rgb]{1,0,0}{T4}}.\\
		
	\end{enumerate}	


%---------------------------------------------------------------------------------
\section*{Results}
Water level, water depths, discharge and velocity results of sub-models \textbf{\textcolor[rgb]{1,0,0}{T1}} to \textbf{\textcolor[rgb]{1,0,0}{T4}} ran on SOBEK3(dflow1d) and DFlowFM are given in \autoref{fig:10404:WaterlevelsAtConNodes} and \autoref{fig:10404:VelocitiesAtU-points}. 

\begin{figure}[H]
    \centering
    \includegraphics*[width=0.85\textwidth]{figures/e02_f104_c04_Waterlevels_at_ConNodes.jpg}
    \caption{Water levels and water depths at (Connection) Nodes}
    \label{fig:10404:WaterlevelsAtConNodes}
\end{figure} 

\begin{figure}[H]
    \centering
    \includegraphics*[width=0.85\textwidth]{figures/e02_f104_c04_Velocities_at_u_points.jpg}
    \caption{Discharges and velocity at u-points directly adjacent to (connection) nodes}
    \label{fig:10404:VelocitiesAtU-points}
\end{figure}

%---------------------------------------------------------------------------------
\section*{Analysis of results}
The results are analyzed through (1) a comparison to the results of the identical Sobek3 Junction Advection test case and; (2) a comparison of the results for sub-models with (i.e. \textbf{\textcolor[rgb]{1,0,0}{T3}} and \textbf{\textcolor[rgb]{1,0,0}{T4}}) and without (i.e. \textbf{\textcolor[rgb]{1,0,0}{T1}} and \textbf{\textcolor[rgb]{1,0,0}{T2}}) junction-advection applied.

From a comparison to the results of the identical Sobek3 Junction-Advection test case, the following can be observed:
\begin{enumerate}

	\item The water depth and water levels at the (connection nodes) are consistently higher for the DFlowFM test case than the SOBEK3 test case. As discharge is maintained in most sub-models (except in sub-model \textbf{\textcolor[rgb]{1,0,0}{T4}}), water velocities are thus also lower around the connection nodes for the DFlowFM test case. This difference is not caused by the application of junction-advection as it is also observed for sub-models \textbf{\textcolor[rgb]{1,0,0}{T1}} and \textbf{\textcolor[rgb]{1,0,0}{T2}} which do not have junction-advection applied. The investigation of this difference therefore falls outside the scope of this test case but should be further investigated. \newline
	
	\item The discharges after the connection node in sub-model \textbf{\textcolor[rgb]{1,0,0}{T4}}are different for the DFlowFM and SOBEK3 case. This is due to some difference in the distribution of flow after the connection node when the cross-sections of the connecting branches vary. The cross section width of branches B2 and B3 are 50m and 450m respectively in sub-model \textbf{\textcolor[rgb]{1,0,0}{T4}}. A simple distribution calculation of the total discharge (790.57 $m^{3}$/s) flowing into the connection node based on a discharge per unit width suggest a distribution in line with the SOBEK3 results. In the DFlowFM case, a larger portion of flow is diverted to B2 causing higher velocities downstream of the connection node. As this difference is not attributed the incorrect application of junction advection, the investigation of this difference falls outside the scope of this test case but should be further investigated. \newline
	
	\item Due to the difference in the distribution of flow after a bifurcation to  branches with varying cross-sectional widths (discussed in point 2 above); the velocities in branches B2 and B3 after(downstream) the connection node are not identical in sub-model \textbf{\textcolor[rgb]{1,0,0}{T4}} for the DFlowFM test case. While they are identical in the SOBEK3 test case. \newline

\end{enumerate}
	
From a comparison of the results for sub-models with (i.e. \textbf{\textcolor[rgb]{1,0,0}{T3}} and \textbf{\textcolor[rgb]{1,0,0}{T4}})and without (i.e. \textbf{\textcolor[rgb]{1,0,0}{T1}} and \textbf{\textcolor[rgb]{1,0,0}{T2}}) junction-advection applied, the following can be observed:

\begin{enumerate}
		
		\item Water depths and levels increase when the straight channel is converted to a bifurcation. This is evident when comparing the results of the straight line sub-models \textbf{\textcolor[rgb]{1,0,0}{T1}} and \textbf{\textcolor[rgb]{1,0,0}{T2}} to bifurcated sub-models \textbf{\textcolor[rgb]{1,0,0}{T3}} and \textbf{\textcolor[rgb]{1,0,0}{T4}}. This increase is in line with the increase in water depth observed for a single channel system with a directional change at that location (see \textbf{f115}\_numerical\_aspects     \_\textbf{c16}\_curved-channel-junction-advection-reference) and is not caused by an incorrect application of junction advection. The investigation of this difference therefore falls outside the scope of this test case but should be further investigated. \newline

	\item In line with the observed increase in the water depth and level in sub-models with a bifurcation as discussed in the point 1 above; the velocities before and after the connection nodes in sub-models \textbf{\textcolor[rgb]{1,0,0}{T3}} and \textbf{\textcolor[rgb]{1,0,0}{T4}} are also reduced. \newline
	
\end{enumerate}

%---------------------------------------------------------------------------------
\subsection{Junction advection} \label{10404:description of junction advection}
Junction advection (or accounting for the advection term ($u \pdiff{u}{x}$) near connection nodes) is incorporated in the dflowFM1d (cf\_dll.dll) computational core in accordance with \citet{Borsboom2017}.

Junction advection is not applied at boundary nodes. For connection nodes yields that junction advection is only applied if more than two branches are connected at a connection node. The reason, herefore, are:
\begin{itemize}
	\item For a branch connected at a boundary node, the advection term ($u \pdiff{u}{x}$) near the boundary node is determined by taking the velocity ($u$) equal to zero at the boundary node.
	\item For connection node having only one branch connected to it, the advection term ($u \pdiff{u}{x}$) near the connection node is determined by taking the velocity ($u$) equal to zero at the connection node.	
	\item For two connected branches, the connection node is just treated as an ordinary water-level point. In other words the advection term ($u \pdiff{u}{x}$) is computed in the same manner as for any water-level point, located in a branch at sufficient distance of a connection node or boundary node.
\end{itemize} 

For connection nodes having \underline{more than two} connected branches, junction advection means that the advection term ($u \pdiff{u}{x}$) is taken into account in solving the St. Venant momentum equation between a connection node (or junction) and the first adjacent water-level-point, located on a particular connected branch.

Before junction advection was implemented (i.e. no junction advection was applied) in the 'dflow1d' computational core following was done. For connection nodes having more than two connected branches, the advection term ($u \pdiff{u}{x}$) directly adjacent to such connection node was set equal to zero in solving the St. Venant momentum equation. 

\Note Junction advection does \textbf{\textcolor[rgb]{1,0,0}{not}} account for (possible) hydraulic head-losses (or \textbf{\textcolor[rgb]{1,0,0}{junction-losses}}) at a connection node. The reason being that junction advection does \underline{not} take into account the angles (in the horizontal and/or vertical plane) between the various branches, that are connected at a connection node (or junction). 

%---------------------------------------------------------------------------------
\printrefsegment
